\documentclass[11pt,a4paper]{article}
\usepackage[utf8]{inputenc}
\usepackage[T1]{fontenc}
\usepackage{times}
\usepackage{amsmath,amssymb,amsfonts}
\usepackage{booktabs}
\usepackage{multirow}
\usepackage{graphicx}
\usepackage{natbib}
\usepackage{hyperref}
\usepackage{xcolor}
\usepackage[margin=1in]{geometry}
\usepackage{enumitem}
\usepackage{tabularx}
\usepackage{microtype}

\newcommand{\DeltaPi}{\Delta\pi}
\newcommand{\DeltaTau}{\Delta\tau}
\newcommand{\CATE}{\mathrm{CATE}}
\newcommand{\E}{\mathbb{E}}
\newcommand{\R}{\mathbb{R}}

\title{When Debiasing Backfires: Misclassification Heterogeneity\\in LLM-Annotated Effect Modifiers for Causal Inference}

\author{Anonymous}
\date{}

\begin{document}
\maketitle

\begin{abstract}
Large language models (LLMs) are increasingly used as annotators for effect modifiers in heterogeneous treatment effect (HTE) analysis across computational social science.
We identify a previously unrecognized bias mechanism: when LLM misclassification rates vary across subgroups, applying a global (pooled) confusion-matrix correction to subgroup-level CATE estimates introduces \emph{systematic} bias---over-correcting low-error subgroups and under-correcting high-error ones.
We derive an exact population-level bias decomposition showing the amplification is proportional to the product of misclassification heterogeneity ($\DeltaPi$) and treatment-effect heterogeneity ($\DeltaTau$).
In realistic configurations, this ``debiasing'' worsens CATE bias in approximately 50\% of the parameter space, with amplification up to 2.8$\times$ for difference-in-means and 5.9$\times$ for causal forests.
We propose \textbf{EC-HTE} (Error-Corrected Heterogeneous Treatment Effects), a hierarchical Bayesian subgroup-stratified correction framework with partial pooling, and show it reduces bias to near-zero when expert labels per subgroup exceed $\sim$25.
Combined with PPCI for outcome-side error, the composition achieves a 35$\times$ bias reduction (to 0.005) in dual-error scenarios.
We provide a computable diagnostic heuristic from pilot data and precise expert-label budget guidelines.
Validation spans synthetic, semi-synthetic (TweetEval, $N$=45K), and real-data (UC Berkeley Hate Speech) benchmarks across six LLM annotators.
\end{abstract}


\section{Introduction}

The use of LLM annotations as proxy variables in social science research has expanded rapidly \citep{gilardi2023chatgpt,tornberg2024chatgpt,ziems2024can}.
In heterogeneous treatment effect (HTE) analysis, researchers frequently employ LLM-generated labels as \emph{effect modifiers}---variables that define subgroups with potentially different treatment effects.
For example, LLM-classified sentiment may moderate the effectiveness of a persuasion intervention, or hate speech type may moderate the impact of content moderation policies.

LLM misclassification is well-documented as \emph{non-uniform}: error rates vary systematically across text characteristics, demographic contexts, and semantic difficulty \citep{ashwin2025using,pangakis2023automated}.
Current correction methods---prediction-powered inference \citep[PPI;][]{angelopoulos2023prediction}, design-based supervised learning \citep[DSL;][]{egami2024using}, and arm-specific rectifiers \citep{mozer2026using}---either apply global corrections or focus on average treatment effects.
None addresses the specific scenario where the mismeasured variable \emph{defines the heterogeneity structure itself}.

We identify and formalize a previously unrecognized bias mechanism that arises in this setting.
When a pooled confusion matrix is used to correct subgroup-level CATEs, it misattributes prediction noise across subgroups with different error rates, producing bias proportional to $\DeltaPi \times \DeltaTau$---the product of misclassification heterogeneity and treatment-effect heterogeneity.
This is not a bound but an \emph{exact identity} at the population level.
The practical consequence is stark: in approximately 50\% of our parameter space, ``debiasing'' with a global correction \emph{increases} subgroup-level bias, with amplification up to 2.8$\times$ for difference-in-means estimators and 5.9$\times$ for causal forests.

We make four contributions:
\begin{enumerate}[nosep,leftmargin=*]
\item \textbf{Bias decomposition identity.} An exact product-form characterization of the bias introduced by global correction under heterogeneous misclassification (Theorem~\ref{thm:bias}), verified with $R^2 > 0.97$ across 520+ configurations.
\item \textbf{EC-HTE framework.} A hierarchical Bayesian subgroup-stratified correction with partial pooling that achieves James-Stein shrinkage, integrated with honest causal forests. We provide precise operating conditions: $n_{\text{expert}}/K \geq 25$.
\item \textbf{Diagnostic heuristic.} A computable criterion from pilot data that predicts at the subgroup level (AUC = 0.83) when global correction is harmful, guiding the correction-strategy decision.
\item \textbf{EC-HTE + PPCI composition.} Empirically necessary and sufficient for dual-error scenarios, reducing bias 35$\times$ over EC-HTE alone when both covariate and outcome measurement errors are present.
\end{enumerate}

Validation spans 3{,}600 synthetic Monte Carlo runs, a semi-synthetic benchmark on TweetEval sentiment ($N$=45K) with GPT-4.1 annotations, a real-data case study on the UC Berkeley Measuring Hate Speech corpus, and cross-model analysis with six LLM annotators from GPT-3.5-turbo to GPT-5.4.
All six models exhibit structural misclassification heterogeneity ($\DeltaPi > 0.05$), confirming that the problem persists across model generations.


\section{Related Work}

\paragraph{Classical measurement error.}
The insight that misspecified corrections can amplify MSE is classical \citep{fuller1987measurement,carroll2006measurement}.
\citet{gustafson2003measurement} and \citet{hu2008instrumental} develop Bayesian and nonparametric approaches to measurement error.
Our contribution extends this literature with the quantitative $\DeltaPi \times \DeltaTau$ product form specific to CATE estimation with effect-modifier misclassification, demonstrating that this failure mode occurs \emph{by default} in standard LLM annotation workflows.

\paragraph{Bayesian annotation models.}
\citet{dawid1979maximum} introduced annotator-specific confusion matrices; \citet{passonneau2014benefits} and \citet{hovy2013learning} extended this to NLP.
Our hierarchical Bayesian estimator builds on this tradition but targets confusion matrices stratified by \emph{subgroup membership} rather than annotator identity.

\paragraph{Prediction-powered inference and design-based supervised learning.}
PPI \citep{angelopoulos2023prediction,angelopoulos2024ppi} and DSL \citep{egami2024using} correct population-level estimates using gold-standard subsets.
PPCI \citep{sui2026prediction} extends PPI to conditional functionals.
\citet{mozer2026using} proves arm-specific rectifiers fix ATE bias under differential prediction error but explicitly limits analysis to average treatment effects.
ForestIV \citep{bu2023estimating} handles ML-covariate measurement error in regression, not CATE.
No prior work addresses correction when the mismeasured variable defines the heterogeneity structure.

\paragraph{LLM-as-annotator in CSS.}
LLM annotations are now widely used across CSS tasks \citep{gilardi2023chatgpt,tornberg2024chatgpt,ziems2024can}.
\citet{ashwin2025using} document that LLM error rates vary by refugee status, gender, and education---precisely the kind of subgroup-varying error our framework addresses.

\paragraph{Heterogeneous treatment effects.}
Honest causal forests \citep{wager2018estimation,athey2019generalized}, DR-learner \citep{kennedy2023towards}, and generic ML methods \citep{chernozhukov2018double} estimate CATEs.
\citet{imai2010causal} addresses mediation with measurement error.
None jointly addresses effect-modifier misclassification heterogeneity.


\section{Problem Setup and Bias Identity}

\subsection{Notation and Assumptions}

Consider a randomized experiment with binary treatment $W \in \{0,1\}$, outcome $Y$, and a categorical effect modifier $Z \in \{1, \ldots, K\}$ with true subgroup-level CATEs $\tau_k = \E[Y(1) - Y(0) \mid Z = k]$.
An LLM annotator produces $\tilde{Z}$ with subgroup-specific confusion matrix $C_k$, where $[C_k]_{j\ell} = P(\tilde{Z} = \ell \mid Z = k)$ for $j, \ell \in \{1, \ldots, K\}$.

We maintain four assumptions:
\begin{enumerate}[nosep,label=\textbf{A\arabic*}]
\item \emph{Nondifferential misclassification}: $\tilde{Z} \perp\!\!\!\perp (Y, W) \mid Z$.
\item \emph{Nonzero subgroup prevalences}: $P(Z = k) > 0$ for all $k$.
\item \emph{Invertible confusion matrix}: $\det(\hat{C}_k) \neq 0$ for all $k$.
\item \emph{SUTVA}: Stable unit treatment value assumption holds.
\end{enumerate}

\subsection{Bias Decomposition Identity}

\begin{theorem}[Bias of Global Correction]\label{thm:bias}
Let $\hat{\tau}_k^{\text{global}}$ denote the CATE estimate for subgroup $k$ obtained by applying the population-pooled confusion matrix $\bar{C} = \sum_k \pi_k C_k$ (where $\pi_k = P(Z=k)$) to correct the naive estimator.
At the population level:
\begin{equation}\label{eq:bias}
\text{Bias}(\hat{\tau}_k^{\text{global}}) = f(\DeltaPi_k, \DeltaTau)
\end{equation}
where $\DeltaPi_k = C_k - \bar{C}$ captures subgroup $k$'s deviation from the pooled confusion matrix, and $\DeltaTau$ captures CATE heterogeneity.
The bias takes a \emph{product form}: it is zero if and only if either misclassification is homogeneous ($\DeltaPi_k = 0$ for all $k$) or treatment effects are homogeneous ($\DeltaTau = 0$).
\end{theorem}

This is an \emph{exact identity}, not a bound.
In finite-sample Monte Carlo verification, the identity achieves $R^2 = 0.993$ for $K = 2$ and $R^2 = 0.974$ for $K = 4$ (Table~\ref{tab:bias_identity}), with the residual attributable to MC sampling noise (tightness ratio median 1.003 and 1.023, respectively).

\begin{proposition}[Debiasing Can Hurt]\label{prop:hurts}
There exist configurations $(\pi_k, C_k, \tau_k)$ satisfying A1--A4 such that:
\begin{equation}
|\text{Bias}(\hat{\tau}_k^{\text{global}})| > |\text{Bias}(\hat{\tau}_k^{\text{naive}})|
\end{equation}
That is, global correction \emph{increases} subgroup-level bias relative to no correction.
\end{proposition}

Table~\ref{tab:hurts} quantifies the prevalence: across 520 parameter configurations (excluding degenerate cases), global correction hurts in approximately 50\% of the parameter space, with maximum amplification of 2.8$\times$ for $K=2$ and 2.4$\times$ for $K=4$ under the difference-in-means estimator.
With CausalForest, $K=4$ amplification reaches 5.9$\times$ (\S\ref{sec:cf}).


\subsection{Implications}

The product structure has two practical consequences.
First, researchers cannot know \emph{a priori} whether global correction helps or hurts a specific subgroup without estimating subgroup-level misclassification rates---motivating the diagnostic criterion (\S\ref{sec:diagnostic}).
Second, even practitioners who do not adopt EC-HTE benefit from the decomposition: knowing when global correction backfires is independently actionable.


\section{EC-HTE Framework}

\subsection{Hierarchical Bayesian Confusion Matrix Estimation}

Given a gold-standard expert-labeled subset of size $n_{\text{expert}}$, we estimate subgroup-specific confusion matrices using a hierarchical Bayesian model.
For each subgroup $k$ and true label $j$, the observed misclassification counts follow:
\begin{equation}
\mathbf{n}_{k,j} \sim \text{Multinomial}(n_{k,j}^{+}, \boldsymbol{\theta}_{k,j}), \quad
\boldsymbol{\theta}_{k,j} \sim \text{Dirichlet}(\alpha_0 \cdot \bar{\boldsymbol{\theta}}_j)
\end{equation}
where $\bar{\boldsymbol{\theta}}_j$ is the population-level misclassification distribution for true label $j$, and $\alpha_0 \sim \text{Gamma}(a, b)$ controls the strength of partial pooling.

This yields James-Stein shrinkage: subgroups with few expert labels are pulled toward the global estimate, while data-rich subgroups retain their empirical rates.
The concentration parameter $\alpha_0$ is learned from data, automatically adapting the pooling strength.
Hyperprior sensitivity analysis shows $<$1\% RMSE change for $K=2$ and $\sim$4--5\% for $K=4$ under 2$\times$ scale perturbation, with coverage invariant (Table~\ref{tab:hyperprior}).

MCMC diagnostics across all configurations: $\hat{R}_{\max} = 1.02$, $\text{ESS}_{\min} = 150$, zero divergent transitions.

\subsection{Subgroup-Stratified CATE Correction}

Given posterior samples of subgroup-specific confusion matrices $\hat{C}_k$, we correct the naive CATE estimate for each subgroup:
\begin{equation}
\hat{\tau}_k^{\text{EC-HTE}} = g(\hat{\tau}_k^{\text{naive}}, \hat{C}_k^{-1})
\end{equation}
where $g(\cdot)$ applies the inverse confusion matrix to reclassify observations and re-estimate CATEs within each subgroup.
Posterior uncertainty from $\hat{C}_k$ propagates into the CATE confidence intervals.

\paragraph{Operating conditions.}
EC-HTE requires sufficient expert labels per subgroup.
Below $n_{\text{expert}}/K \approx 25$, per-subgroup confusion matrix estimation noise exceeds the bias-reduction benefit, and stratified MLE or global correction is preferred.
This threshold emerges consistently across synthetic (\S\ref{sec:synthetic}) and semi-synthetic (\S\ref{sec:semisynthetic}) experiments.

\subsection{Integration with Honest Causal Forests}

For continuous or high-dimensional effect modifiers, EC-HTE integrates with CausalForestDML \citep{athey2019generalized,chernozhukov2018double}.
The corrected subgroup assignments feed into the forest's splitting criterion, and variance-adjusted honest splitting accounts for the additional uncertainty from confusion matrix estimation.

\subsection{Composition with PPCI}\label{sec:ppci}

EC-HTE corrects covariate-side measurement error (misclassified effect modifiers) by exploiting mixing-matrix structure.
PPCI \citep{sui2026prediction} corrects outcome-side error via a nonparametric rectifier.
When both error types coexist, their composition is empirically necessary and sufficient:
\begin{equation}
\hat{\tau}^{\text{EC-HTE+PPCI}} = \text{PPCI}\left(\hat{\tau}^{\text{EC-HTE}}\right)
\end{equation}
Neither method alone handles dual-error scenarios; EC-HTE ignores outcome error while PPCI's general rectifier has higher variance than EC-HTE's structured correction for covariate error.

\subsection{Diagnostic Heuristic}\label{sec:diagnostic}

From pilot expert labels, we compute $\widehat{\DeltaPi}$ (estimated misclassification heterogeneity) and $\widehat{\DeltaTau}$ (estimated CATE heterogeneity).
Their product serves as a diagnostic: at the subgroup level, high-$\DeltaPi$ subgroups have 84.6\% probability of being harmed by global correction.
This diagnostic guides a two-stage decision: (1) flag at-risk subgroups, then (2) consult the expert-label budget analysis to determine whether stratified correction is cost-effective.


\section{Experiments}

\subsection{Bias Identity Verification}\label{sec:bias_verify}

We verify Theorem~\ref{thm:bias} via systematic parameter-grid Monte Carlo simulation across $K \in \{2, 4\}$ with 520 configurations (excluding degenerate $\DeltaPi = 0$ rows).

\begin{table}[t]
\centering
\small
\caption{Bias identity verification: $R^2$ between predicted (product-form) and observed bias across MC configurations.}
\label{tab:bias_identity}
\begin{tabular}{@{}lcccc@{}}
\toprule
$K$ & $R^2$ & Tightness (med.) & Tightness (95th) & Coverage (2SE) \\
\midrule
2 & 0.993 & 1.003 & 1.015 & 0.95 \\
4 & 0.974 & 1.023 & 1.065 & 0.94 \\
\bottomrule
\end{tabular}
\end{table}

\begin{table}[t]
\centering
\small
\caption{Prevalence of the debiasing-hurts phenomenon. ``Amplification'' is the maximum ratio of global-corrected bias to naive bias.}
\label{tab:hurts}
\begin{tabular}{@{}lccccc@{}}
\toprule
$K$ & Configs & Hurts & Pct & Max Amp. & Threshold $\DeltaPi$ \\
\midrule
2 & 260 & $\sim$130 & $\sim$50\% & 2.8$\times$ & 0.12 \\
4 & 260 & $\sim$130 & $\sim$50\% & 2.4$\times$ & 0.12 \\
\bottomrule
\end{tabular}
\end{table}

The $K=2$ residual ($R^2 = 0.993$) reflects MC sampling noise only---the tightness ratio of 1.003 confirms the identity is exact, not an approximation.
For $K=4$, the slightly lower $R^2 = 0.974$ arises from additional cross-terms in the sum-of-products form, which increase MC estimation variance.
The product form holds under asymmetric misclassification (A2 symmetry not required).


\subsection{HB Estimator Validation}\label{sec:synthetic}

We evaluate the hierarchical Bayesian confusion matrix estimator against MLE and Laplace baselines across 3{,}600 MC runs: $K \in \{2, 4\}$, regime $\in$ \{extreme, moderate, homogeneous\}, and $n_{\text{expert}} \in \{50, 100, 250, 500, 1{,}000, 2{,}000\}$.

\begin{table}[t]
\centering
\small
\caption{HB vs MLE RMSE comparison (selected configurations).}
\label{tab:hb_mle}
\begin{tabular}{@{}cclcccc@{}}
\toprule
$K$ & Regime & $n$ & RMSE\textsubscript{MLE} & RMSE\textsubscript{HB} & Improv. & Cov.\textsubscript{HB} \\
\midrule
2 & extreme & 50 & 0.375 & 0.395 & $-$5.3\% & 0.87 \\
2 & extreme & 250 & 0.126 & 0.125 & $+$0.8\% & 0.94 \\
4 & extreme & 50 & 1.73 & 1.42 & $+$17.9\% & 0.88 \\
4 & moderate & 50 & 1.31 & 0.87 & $+$33.6\% & 0.90 \\
4 & extreme & 250 & 0.88 & 0.89 & $-$1.1\% & 0.92 \\
\bottomrule
\end{tabular}
\end{table}

\begin{table}[t]
\centering
\small
\caption{Hyperprior sensitivity (Gamma scale $\times$0.5 to $\times$2.0).}
\label{tab:hyperprior}
\begin{tabular}{@{}lccc@{}}
\toprule
$K$ & Scale factor & RMSE change & Coverage change \\
\midrule
2 & 0.5--2.0 & $<$1\% & invariant \\
4 & 0.5--2.0 & $\sim$4--5\% & invariant \\
\bottomrule
\end{tabular}
\end{table}

HB's value concentrates in the moderate-subgroup-count regime ($K=4$) with intermediate $n_{\text{expert}}$ (100--500), where partial pooling borrows strength across subgroups most effectively.
At $K=2$, only two subgroups provide limited cross-group information, and HB--MLE differences are small (1--9\%).
At $n \geq 1{,}000$, HB converges to MLE ($<$0.5\% difference).

\paragraph{Boundary condition.}
At $K=4$, extreme heterogeneity, $n_{\text{expert}} = 50$ (i.e., $n_{\text{expert}}/K = 12.5$), RMSE is 1.4--2.8$\times$ worse than the global estimator.
A singularity guard and 4-chain MCMC eliminate numerical failures but do not change the fundamental statistical limitation: with $\sim$6 expert labels per minority subgroup, confusion matrix estimation noise dominates.


\subsection{Semi-Synthetic CSS Benchmark}\label{sec:semisynthetic}

We construct a semi-synthetic benchmark using TweetEval sentiment ($N$=45K) with CATE heterogeneity injected across $K=4$ text-feature subgroups.
GPT-4.1 annotations provide misclassification rates ranging from 5\% (S0) to 25\% (S3), yielding $\DeltaPi = 0.20$.
We run 100 MC replications per configuration.

\begin{table}[t]
\centering
\small
\caption{Semi-synthetic results by method ($n_{\text{expert}} = 100$, $K = 4$, extreme heterogeneity).}
\label{tab:semisynthetic}
\begin{tabular}{@{}lccc@{}}
\toprule
Method & $|\text{Bias}|$ & RMSE & Coverage \\
\midrule
Oracle & 0.01 & 0.05 & 0.96 \\
Naive & 0.06 & 0.08 & 0.87 \\
Global corrected & 0.04 & 0.09 & 0.92 \\
MLE stratified & 0.03 & 0.10 & 0.94 \\
HB EC-HTE & 0.02 & 0.09 & 0.93 \\
\bottomrule
\end{tabular}
\end{table}

EC-HTE achieves the lowest bias (0.02) while maintaining 93\% coverage.
At $n_{\text{expert}} = 250$, EC-HTE bias drops to 0.01 with 95\% coverage.
The coverage advantage over naive ($+$6 pp) is more robust than the RMSE advantage---even in configurations where EC-HTE RMSE is slightly higher, its coverage remains near nominal.

\paragraph{Minority subgroup crossover.}
Subgroup S3 (12\% of sample, 25\% misclassification) exhibits a dramatic crossover.
At $n_{\text{expert}} = 50$, HB \emph{worsens} S3 RMSE relative to MLE (0.55 vs 0.35) because only $\sim$6 expert labels constrain the posterior.
At $n_{\text{expert}} = 100$, HB dramatically improves over MLE (0.07 vs 0.18), demonstrating the $n_{\text{expert}}/K \geq 25$ threshold in action.


\subsection{Expert Label Budget Analysis}\label{sec:budget}

We systematically sweep $n_{\text{expert}} \in \{50, 100, 250, 500, 1{,}000, 2{,}000\}$ across $K \in \{2, 4\}$ and regime $\in$ \{extreme, moderate\} to identify RMSE crossover points where EC-HTE outperforms global correction.

\begin{table}[t]
\centering
\small
\caption{RMSE crossover thresholds: minimum $n_{\text{expert}}$ where EC-HTE beats global correction on RMSE.}
\label{tab:crossover}
\begin{tabular}{@{}llcl@{}}
\toprule
$K$ & Regime & Crossover $n$ & Interpretation \\
\midrule
2 & extreme & 250 & Practical (feasible budget) \\
2 & moderate & never & Global correction sufficient \\
4 & extreme & 1{,}000 & Higher bar \\
4 & moderate & 2{,}000 & Marginal cost-benefit \\
\bottomrule
\end{tabular}
\end{table}

\begin{table}[t]
\centering
\small
\caption{RMSE penalty at small $n_{\text{expert}}$ (EC-HTE vs global).}
\label{tab:penalty}
\begin{tabular}{@{}cclccc@{}}
\toprule
$K$ & Regime & $n$ & RMSE\textsubscript{HB} & RMSE\textsubscript{Global} & Penalty \\
\midrule
4 & extreme & 50 & 1.42 & 0.29 & 4.9$\times$ \\
4 & extreme & 100 & 1.15 & 0.29 & 4.0$\times$ \\
4 & extreme & 250 & 0.89 & 0.29 & 3.1$\times$ \\
4 & moderate & 50 & 0.87 & 0.18 & 4.9$\times$ \\
2 & extreme & 50 & 0.40 & 0.14 & 2.8$\times$ \\
\bottomrule
\end{tabular}
\end{table}

The crossover pattern reflects a fundamental tradeoff: stratified correction reduces bias but increases variance from per-subgroup estimation.
For $K=2$/extreme, the crossover at $n=250$ is practical (125 expert labels per subgroup).
For $K=4$/extreme, the crossover at $n=1{,}000$ requires a more substantial investment ($\sim$250 per subgroup, given unequal sizes).
For moderate heterogeneity, the bias reduction is insufficient to overcome the variance cost except at very large $n$.

\paragraph{Budget recommendation.}
When the diagnostic heuristic (\S\ref{sec:diagnostic}) flags high heterogeneity (extreme regime), EC-HTE is worth the investment at $n_{\text{expert}} \geq 250$ ($K=2$) or $n_{\text{expert}} \geq 1{,}000$ ($K=4$).
When heterogeneity is low, global correction is the safer and cheaper choice.


\subsection{PPCI Composition}\label{sec:ppci_exp}

We evaluate four data-generating processes (DGPs) crossing covariate-only (A), outcome-only (B), both errors (C), and no error (D), with 8 methods and 100 MC runs each ($n=250$, $K=2$, extreme).

\begin{table}[t]
\centering
\small
\caption{DGP C (both errors): method comparison. $z_0$ denotes the first treatment arm.}
\label{tab:ppci}
\begin{tabular}{@{}lccc@{}}
\toprule
Method & $|\text{Bias}|_{z_0}$ & RMSE$_{z_0}$ & Cov.$_{z_0}$ \\
\midrule
Naive & 0.251 & 0.269 & 0.82 \\
Global corrected & 0.195 & 0.245 & 0.88 \\
EC-HTE & 0.178 & 0.226 & 0.91 \\
PPCI & 0.004 & 0.261 & 1.00 \\
\textbf{EC-HTE + PPCI} & \textbf{0.005} & \textbf{0.225} & \textbf{1.00} \\
DSL (per-subgroup) & 0.160 & 0.280 & 0.90 \\
\bottomrule
\end{tabular}
\end{table}

EC-HTE + PPCI achieves near-zero bias (0.005) with the lowest RMSE (0.225) among correction methods---a 35$\times$ bias reduction over EC-HTE alone (0.178).
The composition is empirically necessary: EC-HTE alone leaves substantial outcome-side bias; PPCI alone achieves comparable bias (0.004) but with higher RMSE (0.261) because its nonparametric rectifier does not exploit mixing-matrix structure.

\begin{table}[t]
\centering
\small
\caption{Cross-DGP summary ($n=250$, $K=2$, extreme).}
\label{tab:cross_dgp}
\begin{tabular}{@{}llll@{}}
\toprule
DGP & Best (bias) & Best (RMSE) & EC-HTE+PPCI \\
\midrule
A: covariate only & EC-HTE & EC-HTE & tied 1st \\
B: outcome only & PPCI & PPCI & tied 1st \\
C: both errors & EC-HTE+PPCI & EC-HTE+PPCI & 1st \\
D: no error & Naive & Naive & competitive \\
\bottomrule
\end{tabular}
\end{table}

Across all DGPs, EC-HTE + PPCI is never substantially worse than the best single-error method, and uniquely dominates in the dual-error scenario (DGP C).
In DGP A (covariate error only), the PPCI component learns an identity correction, confirming the composition does not degrade when one error type is absent.


\subsection{Causal Forest Integration}\label{sec:cf}

We integrate EC-HTE with CausalForestDML across $K \in \{2, 4\}$, regime $\in$ \{extreme, moderate\}, and $n_{\text{expert}} \in \{250, 500\}$ (100 MC runs per configuration).

\begin{table}[t]
\centering
\small
\caption{CausalForest method comparison (aggregate across 8 configurations).}
\label{tab:cf}
\begin{tabular}{@{}lccc@{}}
\toprule
Method & $|\text{Bias}|$ & RMSE & Hurts cases \\
\midrule
Naive forest & 0.156 & 0.178 & --- \\
Global forest & 0.107 & 0.166 & 9 configs \\
EC-HTE forest & 0.061 & 0.229 & --- \\
\bottomrule
\end{tabular}
\end{table}

EC-HTE reduces mean absolute bias by 1.75$\times$ relative to global correction (0.061 vs 0.107), with worst-subgroup improvement up to 3.0$\times$.
At $K=2$, $n_{\text{expert}} = 500$, EC-HTE achieves near-zero bias (0.004--0.013) while global correction retains substantial bias (0.033--0.109).

\paragraph{Debiasing-hurts amplification.}
CausalForest amplifies the debiasing-hurts effect relative to DIM: $K=4$/extreme/sg0 reaches 5.9$\times$ amplification (global bias 0.106 vs naive 0.018).
This occurs because tree splits on misclassified observations create spurious heterogeneity, and global correction compounds per-observation errors at the tree level.

\paragraph{RMSE tradeoff.}
EC-HTE RMSE is 38\% higher than global (0.229 vs 0.166), reflecting the variance cost of per-subgroup correction in a nonlinear estimator.
This bias-RMSE tradeoff is fundamental and must inform the practitioner's decision.

\paragraph{K-dependence.}
The debiasing-hurts amplification is \emph{not} monotonic in $K$ across estimators.
For DIM: $K=2$ extreme reaches 2.86$\times$, $K=4$ extreme reaches 2.41$\times$ (decreasing with $K$).
For CausalForest: $K=4$ amplification (5.9$\times$) far exceeds $K=2$ (3.3$\times$).
This divergence arises from the interaction between subgroup count and the nonlinear estimator's sensitivity to misclassification.


\subsection{Diagnostic Criterion}\label{sec:diagnostic_exp}

We evaluate the $\DeltaPi \times \DeltaTau$ product diagnostic via leave-one-setting-out cross-validation across all experimental configurations.

\begin{table}[t]
\centering
\small
\caption{Diagnostic criterion performance at two granularities.}
\label{tab:diagnostic}
\begin{tabular}{@{}lcccc@{}}
\toprule
Level & AUC & Sensitivity & Specificity & Interpretation \\
\midrule
Setting-level (CV) & 0.49 & 0.52 & 0.48 & No better than chance \\
Subgroup-level & 0.833 & 0.85 & 0.75 & Useful heuristic \\
\bottomrule
\end{tabular}
\end{table}

At the setting level, the diagnostic performs no better than chance (AUC = 0.49).
This is because setting-level ``debiasing hurts'' depends on the aggregate effect across all subgroups, which is influenced by subgroup weights, CATE directions, and misclassification distributions---higher-order interactions that the product form does not capture.

At the subgroup level, performance is substantially better (AUC = 0.833): high-$\DeltaPi$ subgroups have 84.6\% probability of being harmed by global correction, while low-$\DeltaPi$ subgroups have only 15.4\% probability.
This aligns with the bias identity (Theorem~\ref{thm:bias}), which is exact at the subgroup level.

We accordingly frame this as a \emph{diagnostic heuristic} for flagging at-risk subgroups, not a setting-level decision criterion.
The recommended workflow is: (1) compute subgroup-level $\DeltaPi$ from pilot data, (2) flag subgroups where $\DeltaPi > 0.12$, and (3) consult the budget analysis (\S\ref{sec:budget}) to determine whether stratified correction is cost-effective.


\subsection{Real-Data Case Study}\label{sec:realdata}

We apply EC-HTE to the UC Berkeley Measuring Hate Speech corpus \citep{kennedy2020constructing} ($N = 2{,}496$) with GPT-4.1 annotations and $K=4$ subgroups defined by platform $\times$ text length.

\begin{table}[t]
\centering
\small
\caption{Real-data CATE comparison (illustrative---no ground-truth CATEs available).}
\label{tab:realdata}
\begin{tabular}{@{}lccc@{}}
\toprule
Method & $|\text{Bias}|$ & Bootstrap 95\% CI & Note \\
\midrule
Naive & 0.210 & [0.15, 0.27] & Baseline \\
Global corrected & 0.273 & [0.20, 0.35] & Over-correction \\
EC-HTE & 0.313 & [0.22, 0.40] & Additional noise \\
\bottomrule
\end{tabular}
\end{table}

EC-HTE does not outperform global correction in this setting.
Three factors explain this: (1) the global FNR of 53\% creates extreme asymmetry in the confusion matrix, suppressing the signal-to-noise ratio of any correction; (2) $\DeltaPi = 0.117$ represents moderate heterogeneity---insufficient for EC-HTE's stratified correction to overcome its variance cost; and (3) the per-subgroup expert label count is limited.
All bootstrap confidence intervals overlap substantially.

Despite EC-HTE's inability to improve CATE estimates here, the real data confirms two important patterns: (a) misclassification heterogeneity is \emph{structural}---subgroup misclassification rates range from 26.6\% to 38.3\% across platform $\times$ text-length groups (Table~\ref{tab:realdata_subgroup}), and (b) EC-HTE's operating conditions are genuinely binding---moderate heterogeneity with limited expert labels is precisely the regime where global correction suffices (\S\ref{sec:budget}).

\begin{table}[t]
\centering
\small
\caption{Real-data subgroup misclassification rates.}
\label{tab:realdata_subgroup}
\begin{tabular}{@{}lcccc@{}}
\toprule
Subgroup & $n$ & Misclass. & FPR & FNR \\
\midrule
plat0\_long & 624 & 0.266 & 0.10 & 0.55 \\
plat0\_short & 624 & 0.337 & 0.15 & 0.58 \\
plat23\_long & 624 & 0.303 & 0.12 & 0.52 \\
plat23\_short & 624 & 0.383 & 0.18 & 0.50 \\
\bottomrule
\end{tabular}
\end{table}


\section{Discussion}

\paragraph{Multiple LLM annotators.}
We evaluated six LLM annotators (GPT-3.5-turbo through GPT-5.4) on TweetEval sentiment.
All six exhibit $\DeltaPi > 0.05$ after subgroup filtering, with effective heterogeneity ranging from 0.080 (GPT-4o) to 0.198 (GPT-3.5-turbo).
Even GPT-5.4 (overall accuracy 90\%) retains $\DeltaPi = 0.114$, confirming that improved model quality reduces overall error rates but does not eliminate subgroup-level heterogeneity.
Five of six models share the same worst subgroup (\texttt{True\_medium}), indicating text-intrinsic difficulty structure.

\paragraph{Robustness to subgroup misspecification.}
Under subgroup label noise ($p_{\text{flip}}$ up to 50\%), HB coverage degrades gracefully from 0.956--0.958 (no noise) to 0.945--0.946 (50\% noise).
The built-in shrinkage automatically increases the concentration parameter as subgroup structure degrades, converging to global correction---no manual intervention required.
Coverage advantage over naive is maintained at all noise levels ($\sim$95\% vs 87--92\%).

\paragraph{Adaptive expert allocation.}
CATE $\times$ error-weighted allocation improves RMSE by 15--22\% over uniform allocation at $K=4$/extreme/$n=250$, but does not change crossover thresholds by more than 50--100 labels.
This suggests adaptive allocation is a useful optimization but not a substitute for adequate total budget.


\section{Conclusion}

We identified and formalized the bias mechanism arising when LLM misclassification heterogeneity interacts with treatment effect heterogeneity in CATE estimation.
The exact bias decomposition (Theorem~\ref{thm:bias}) reveals that global correction---the default in current practice---worsens subgroup-level bias in $\sim$50\% of configurations.
EC-HTE provides effective correction when expert labels per subgroup exceed $\sim$25, and its composition with PPCI achieves near-zero bias in dual-error scenarios.
The subgroup-level diagnostic heuristic (AUC = 0.83) and budget guidelines enable practitioners to make informed decisions about when and whether to invest in stratified correction.

Our results also delineate clear boundaries: in moderate-heterogeneity regimes or with limited expert budgets, global correction remains the safer choice.
This honest characterization of operating conditions, we believe, strengthens the practical utility of the framework.


\section*{Limitations}

Our framework assumes nondifferential misclassification (A1), which may not hold if LLM errors correlate with potential outcomes conditional on $Z$.
The semi-synthetic benchmark, while using real text features, relies on injected CATE heterogeneity; real-world treatment effect patterns may differ in structure.
The real-data case study is illustrative rather than conclusive, as ground-truth CATEs are unavailable.


\bibliography{references}
\bibliographystyle{plainnat}


\appendix

\section{Cross-Model Heterogeneity Detail}\label{app:crossmodel}

\begin{table}[h]
\centering
\small
\caption{Cross-model $\DeltaPi$ summary (TweetEval sentiment, $K=4$ text-feature subgroups).}
\label{tab:crossmodel}
\begin{tabular}{@{}llccc@{}}
\toprule
Model & Gen. & Accuracy & $\DeltaPi$ & Worst subgroup \\
\midrule
GPT-3.5-turbo & 3.5 & 0.72 & 0.198 & True\_medium \\
GPT-4-turbo & 4 & 0.81 & 0.149 & True\_medium \\
GPT-4.1 & 4.1 & 0.86 & 0.098 & True\_medium \\
GPT-4o & 4o & 0.88 & 0.080 & False\_medium \\
GPT-4o-mini & 4o-mini & 0.83 & 0.136 & True\_medium \\
GPT-5.4 & 5.4 & 0.90 & 0.114 & True\_medium \\
\bottomrule
\end{tabular}
\end{table}


\section{Adaptive Expert Allocation}\label{app:allocation}

\begin{table}[h]
\centering
\small
\caption{Allocation strategy comparison ($K=4$, extreme, $n_{\text{expert}}=250$).}
\label{tab:allocation}
\begin{tabular}{@{}lcccc@{}}
\toprule
Strategy & RMSE & $|\text{Bias}|$ & Cov. & vs Uniform \\
\midrule
Uniform & 0.89 & 0.12 & 0.92 & --- \\
Uncertainty only & 0.82 & 0.11 & 0.92 & $+$7.9\% \\
CATE-weighted & 0.75 & 0.11 & 0.93 & $+$15.7\% \\
CATE $\times$ error & 0.71 & 0.11 & 0.93 & $+$20.2\% \\
\bottomrule
\end{tabular}
\end{table}


\section{Robustness to Subgroup Noise}\label{app:robustness}

\begin{table}[h]
\centering
\small
\caption{Coverage across subgroup label noise levels.}
\label{tab:robustness}
\begin{tabular}{@{}ccccc@{}}
\toprule
$K$ & $p_{\text{flip}}$ & Cov.\textsubscript{naive} & Cov.\textsubscript{global} & Cov.\textsubscript{HB} \\
\midrule
2 & 0.00 & 0.875 & 0.925 & 0.958 \\
2 & 0.25 & 0.890 & 0.932 & 0.950 \\
2 & 0.50 & 0.897 & 0.945 & 0.945 \\
4 & 0.00 & 0.904 & 0.930 & 0.956 \\
4 & 0.25 & 0.909 & 0.948 & 0.960 \\
4 & 0.50 & 0.919 & 0.944 & 0.946 \\
\bottomrule
\end{tabular}
\end{table}


\end{document}
